\chapter[SYSTEM ANALYSIS AND DESIGN]{System Analysis and Design}

This chapter is organized in the same order as TrustFin AI, starting from the browser interface, through the prediction service, to the stored assessment. It describes the architecture, processing order, class and package responsibilities, deployment view and database design. The diagrams show various views of the same workflow, and the descriptions identify parts of the application corresponding to the views.

\section{System Architecture and Workflow}
\subsection{System Overview}
TrustFin AI consists of a React client and a FastAPI backend. The client gathers application information and displays the results, while the back-end does model inference, explanation generation, validation and persistence. There are separate home-loan, personal-loan and legacy pipelines in a model registry. The chosen pipeline provides its preprocessor, classifier, feature function and explainers.

This separation is significant because there are different inputs to the loan products. Property information is part of the home-loan flow, and purpose and card utilisation is part of the personal-loan flow. Shared route logic coordinates the assessment without the frontend needing to know the feature vector encoded in the model.

\subsection{System Architecture}
Figure~\ref{fig:system-architecture} shows the assessment workflow from the user interface to the results and PDF report. The backend selects the home-loan or personal-loan pipeline, prepares the application data and uses the saved model to make a prediction. The models are trained offline and are not retrained for each application.

The project policy check applies the rules for the selected product alongside the model prediction. SHAP and LIME provide explanations of the prediction, and the faithfulness, stability and consistency checks are combined into a Trust Score. A separate fairness check examines sensitivity to a change in Gender and does not contribute to the Trust Score.

The backend returns these results for display and can generate an assessment PDF. Assessment history is stored using SQLite or a configured PostgreSQL database. Keeping the model result, policy check and explanation scores separate helps the user review the different parts of the assessment.

\reportfigure{figures/diagrams/system_architecture_layered.png}{0.99}{18.2}{System architecture}{Architecture of the TrustFin AI assessment workflow.}{fig:system-architecture}

\subsection{Processing Methodology}
Figure~\ref{fig:methodology} links data preparation, model training, loan assessment, explanation and decision support. The project uses generated home-loan and personal-loan records. Financial features, encoding and scaling prepare the data for product-specific models trained with fixed parameters. The models are evaluated on a hold-out split as well as independent synthetic and boundary datasets.

\begin{samepage}
Data preparation and model training are performed offline. For a new application, the fitted preprocessor and model are used to obtain a prediction, and the project policy is checked for the selected product. SHAP and LIME explain the model output. Faithfulness, stability and consistency contribute equally to the Trust Score, while the fairness check is reported separately. The results can be viewed in the dashboard, saved in assessment history and exported as a PDF.
\end{samepage}

\reportfigure{figures/diagrams/methodology_updated.png}{0.99}{8.0}{Processing methodology}{Processing stages from dataset preparation to explanation and decision support.}{fig:methodology}

\subsection{Assessment and Retrieval Paths}
A new evaluation does inference and explanation before it returns its response. History retrieval is different: it reads the stored data and builds the response object. This prevents giving a new calculation explanation as though it were the first one stored in the dossier.

Row validation is performed before the assessment and the rows that are not usable are marked and their reasons are recorded in the batch processing. The same product pipeline is used for usable records. The batch summary then provides processed results and the overall Trust Score. The persistence errors and transaction handling are still significant operational constraints, since a computed response is a different event from a successfully stored response.

\clearpage
\section{Data Flow Diagrams}
Data flow views are concerned with the information that enters, leaves and flows through the system. They support the architecture by presenting the inputs and outputs of the assessment and the records that are stored, rather than each software module as a deployed service.

\subsection{Level 0: Context Diagram}
The applicant or reviewing officer provides the details of the application to TrustFin AI and gets an assessment. The system boundary includes prediction, explanation and presentation of the result. When users click on a saved application or export a PDF form, there are further interactions with the same assessment system.

\reportfigure{figures/diagrams/context_level_dfd.png}{0.99}{10.5}{Level 0 DFD -- Context Diagram}{Context-level data flow between portal users and the TrustFin AI assessment system.}{fig:context-dfd}
\clearpage

\subsection{Level 1: Internal Data Flow}
The Level 1 view is a separation of input processing, product feature engineering, prediction, explanation generation, validation and persistence. Selected product pipeline receives normalised application data. The answer includes model output, financial calculations and explanation diagnostics, and evaluation records and batch summaries are passed to storage. History retrieval reads those records, and the reporting path turns an assessment into a PDF.

\reportfigure{figures/diagrams/level_1_dfd.png}{0.99}{10.5}{Level 1 DFD -- Internal Data Flow}{Data flow through application processing, prediction, explanation, validation, storage and reporting.}{fig:level1-dfd}
\clearpage

\subsection{Conceptual Entity--Relationship Diagram}
Figure~\ref{fig:conceptual-er} presents a conceptual view of the information involved in an assessment. A loan application receives an assessment containing the model result and approval probability. The assessment includes SHAP and LIME explanations and a validation result. The fairness check assesses Gender sensitivity separately from the three explanation checks used in the Trust Score.

A batch upload can group multiple application assessments, while an individual assessment need not belong to a batch. The diagram shows conceptual responsibilities rather than separate database tables for each box. In the current database, assessment information and explanations are stored in prediction records and JSON fields. The optional batch association uses \texttt{batch\_id} and is managed by the application rather than enforced by a database foreign key.

\reportfigure{figures/diagrams/conceptual_er_diagram.png}{0.99}{14.0}{Conceptual ER diagram}{Conceptual entities and relationships in the assessment workflow.}{fig:conceptual-er}
\clearpage

\section{Software Design}
\subsection{Component Responsibilities}
The component view of Figure~\ref{fig:component-diagram} organizes the software into components. Form entry, product navigation, results, history and batch interaction are handled by React components. The FastAPI router coordinates the model, product policy check, explainers, validation functions and persistence layer. The reporting module generates A4 PDF assessment reports.

The model and explanation computation is done in the backend process. The database interface is provided by SQLAlchemy, and SQLite is used locally while PostgreSQL can be configured via the environment. The components exchange typed request and response structures instead of display text being passed between the model and browser in separate formats.

\reportfigure{figures/diagrams/component_diagram_updated.png}{0.96}{12.4}{Component diagram}{Frontend, backend, prediction, policy, explanation, validation, reporting and persistence components.}{fig:component-diagram}

\subsection{Class Design and Data Contracts}
Figure~\ref{fig:class-diagram} shows selected classes and fields used in an assessment. The \texttt{LoanApplicationRequest} carries the submitted details, and the API returns a \texttt{LoanPredictionResponse} containing the model result, explanations, Trust Score and, when available, the policy check.

The response contains SHAP and LIME explanations represented by \texttt{StandardExplanation}, each with feature contributions. The \texttt{TrustEvaluation} model carries the explanation scores and a separate fairness value; fairness is not included in the Trust Score calculation. A saved \texttt{PredictionLog} can retain the response as JSON. Table~\ref{tab:design-source-map} relates these responsibilities to the source components.

\reportfigure{figures/diagrams/class_diagram_source.png}{0.98}{15.4}{Class diagram}{Selected source classes and data contracts used in loan assessment.}{fig:class-diagram}

\begin{table}[htbp]\centering\small
\caption{Design responsibilities and source components}\label{tab:design-source-map}
\begin{tabularx}{\textwidth}{p{3.8cm}X}\toprule
Responsibility & Source component\\\midrule
Application data & \texttt{LoanApplicationRequest} and product normalisation.\\
Prediction response & \texttt{LoanPredictionResponse} and the product pipeline.\\
Local explanations & \texttt{SHAPExplainer}, \texttt{LIMEExplainer} and \texttt{StandardExplanation}.\\
Trust evaluation & Four validation functions and \texttt{TrustEvaluation}.\\
Dashboard & React applicant, summary, audit and insight components.\\
Report generation & ReportLab PDF generation function and API endpoint.\\
Saved assessment & \texttt{PredictionLog} and \texttt{BatchUpload}.\\\bottomrule
\end{tabularx}\end{table}

The diagram shows selected fields rather than every field in the source classes. Assessment and batch records are stored in the persistence schema. Applicant and officer modes select different client presentations; server-side authentication of accounts is not included.

\subsection{Use Case Diagram}
Figure~\ref{fig:use-case-diagram} shows the main interactions available in the TrustFin AI interface. A user can select a loan type, enter application details, view the assessment and policy check, explore an affordability scenario, and preview or download an assessment report.

The bank-officer view also provides detailed explanation and diagnostic checks, batch upload, assessment history and an applicant-view preview. The user and bank-officer actors represent interface modes. They do not imply authenticated accounts or server-side role enforcement in the current application.

\reportfigure{figures/diagrams/use_case_diagram_updated.png}{0.99}{14.5}{Use case diagram}{Use cases available in the applicant and bank-officer interface modes.}{fig:use-case-diagram}

\subsection{Sequence of an Individual Evaluation}
The sequence diagram in figure~\ref{fig:sequence-diagram} traces one application request. The user enters the values via the dashboard and the backend selects the product pipeline. There is normalisation and feature engineering before the call to the classifier. The transformed instance is then sent to SHAP and LIME, and the component scores are generated in the validation stage.

Completed response is sent back to the client for display. The response is stored with the input record using a route and an evaluation is identified using a dossier reference. The optional PDF interaction takes the assessment and applicant data and creates a report. It is an independent action to the prediction of the application.

The financial summary is based on the model output. Note that the current path does not automatically result in an approved model label being generated when the FOIR is above the configured boundary, but rather a warning is generated and an insolvency status is generated. This is the reason that prediction and financial evaluation are kept as separate outputs.

\reportfigure{figures/images/Sequence_Diagram_Updated.png}{0.99}{16.5}{Evaluation sequence}{Sequence of interactions for application evaluation and optional PDF reporting.}{fig:sequence-diagram}

\subsection{Package Organisation}
The package view in Fig.~\ref{fig:package-diagram} bundles the following sub-packages: presentation, application coordination, prediction, explainability, trust evaluation and data management. The source is not deployed as a separate service per package, but rather via frontend components and backend modules.

The router is located at the coordination point. It triggers the pre-processing and model inference of the product, passes the results to the explanation and validation functions, and generates the response. The persistence is managed by the database code, and the exported document is managed by the PDF module. This layout assigns a meaningful location to each part of the project and enables to trace the origin of a value displayed.

The current application has no account repository or authentication module. It's data-management implementation is geared toward evaluation and batch records. The package boundaries are used to organise the assessment workflow and do not give access rights to the roles the client has selected.

\reportfigure{figures/images/Package_Diagram.jpeg}{0.87}{18.2}{Package diagram}{Logical package organisation of the financial decision-support application.}{fig:package-diagram}

\clearpage
\section{Database and API Design}
\subsection{Database Design}
The database includes the applicant summary fields, raw input JSON, model output, solvency values, SHAP and LIME JSON, three explanation scores, a separate fairness value and the complete response. For presentation and lookup a special dossier reference is used.

The upload reference, product, filename, timestamp, counts of records, and summary of flawed rows and information are stored in \texttt{BatchUpload}. Prediction rows can have the same \texttt{batch\_id} value. This establishes an application-level association between a batch and its evaluations, but the SQLAlchemy schema doesn't define a foreign-key constraint for it.

\begin{table}[htbp]\centering\small
\caption{Persistence entities and their purpose}\label{tab:database-entities}
\begin{tabularx}{\textwidth}{p{3.1cm}XX}\toprule
Entity & Identification & Stored information\\\midrule
\texttt{PredictionLog} & Numeric ID and unique dossier reference & Input, model and explanation output, financial summary, component scores and optional batch reference.\\
\texttt{BatchUpload} & Numeric ID and unique upload reference & File/product details, row counts, flawed records and aggregate summary.\\\bottomrule
\end{tabularx}\end{table}

JSON can be used to capture nested explanation information and full responses. History summaries are helpful when using scalar fields but not when decoding each payload. Schema setup utilizes SQLAlchemy table creation and some additive column checks. A versioned database-migration framework is not included.

\subsection{API Boundaries}
The loan router is installed at \texttt{/api/v1/loan}. It provides the following main operations: see Table~\ref{tab:api-endpoints}. All the submitted application, prediction response, stored summaries and PDF request are defined by pydantic schemas. Shared response contract allows the interface and database to have the same assessment structure.

\begin{table}[htbp]\centering\small
\caption{Principal API operations}\label{tab:api-endpoints}
\begin{tabularx}{\textwidth}{p{1.8cm}p{5.0cm}L}\toprule
Method & Relative path & Purpose\\\midrule
POST & \path{/predict} & Evaluate one application.\\
POST & \path{/predict-batch} & Process a CSV upload.\\
POST & \path{/predict-batch-json} & Process JSON batch rows.\\
GET & \path{/evaluations} & List saved evaluations.\\
GET & \texttt{/evaluations/\{id\}} & Retrieve an evaluation.\\
DELETE & \texttt{/evaluations/\{id\}} & Delete an evaluation.\\
GET & \path{/batches} & List saved uploads.\\
GET & \texttt{/batches/\{id\}} & Retrieve a batch.\\
DELETE & \texttt{/batches/\{id\}} & Delete a batch.\\
GET & \path{/batch-sample-csv} & Download sample product input.\\
GET & \texttt{/batch-presets/\{id\}} & Retrieve a demonstration preset.\\
POST & \path{/generate-pdf} & Export the assessment.\\\bottomrule
\end{tabularx}\end{table}

\section{Deployment View}
The deployment view of Figure~\ref{fig:deployment-diagram} is a separation of the browser, application runtime, model responsibilities and database. In local development, the browser will interact with Vite and the FastAPI service. The API loads model artifacts from the file system and does inference and explanation in the Python process.

SQLite provides a local storage place. Or a configured PostgreSQL connection can be used, and a local fallback if the connection fails. The logical AI environment separates the prediction and explanation tasks in the diagram, and the current code does not require a separate model API. Likewise, the security of the transport relies on the configuration and is not defined by a development URL.

\reportfigure{figures/images/deployment-diagram.png}{0.97}{12.7}{Deployment diagram}{Deployment responsibilities for the client, application runtime, model execution and database.}{fig:deployment-diagram}

\par\medskip\noindent
This chapter has been followed up in the application architecture, components, data contracts, processing sequence, package organisation, deployment responsibilities and database. It has also demonstrated how the API can link a new assessment with the history and reporting capabilities.
